% e-TeX's \unless -- previously not implemented at all.
%
% A prefix, not a conditional of its own: it must be followed immediately by
% a boolean \if primitive (resolved through \let, since expl3 renames every
% one of them) and inverts that conditional's result. expl3's
% \reverse_if:N is \let straight to it --
%
%   \tex_let:D \reverse_if:N \tex_unless:D
%
% -- and every negated comparison in the kernel is built on it: \int_compare,
% \dim_compare, \fp_compare, dozens of others. Missing it entirely left
% \reverse_if:N undefined, which cascaded into a runaway a few thousand
% commands later.
\input plain
\def\M#1{\message{<#1>}}
\M{a=\unless\ifnum 1>0 T\else F\fi}
\M{b=\unless\ifnum 1<0 T\else F\fi}
\M{c=\unless\iftrue T\else F\fi}
\M{d=\unless\iffalse T\else F\fi}
\M{e=\unless\ifodd 4 T\else F\fi}
\M{f=\unless\ifodd 3 T\else F\fi}
\M{g=\unless\ifx\undefined\relax T\else F\fi}
\let\samecs\relax
\M{h=\unless\ifx\samecs\relax T\else F\fi}
\M{i=\unless\ifdim 1pt>2pt T\else F\fi}
% nested: unless inside an ordinary conditional's operand
\M{j=\ifnum\unless\ifodd 3 1\else 0\fi=0 T\else F\fi}
% doubled: unless applied where the inner conditional is itself aliased
\let\myifnum\ifnum
\M{k=\unless\myifnum 5>10 T\else F\fi}
\M{l=\ifdefined\reverse_if:N y\else n\fi}
% The one that mattered: \unless prefixes a conditional whose *operand scan*
% opens another conditional. The nested one must not steal the \unless. This
% is the shape \int_compare's loop builds -- \unless\ifnum A < \ifx.. -- where
% a stolen \unless inverted a false test and produced \__int_compare_<=:NNw.
\M{m=\unless\ifnum 5<\ifnum1>0 9\else 1\fi\relax T\else F\fi}
\M{n=\unless\ifnum 5<\ifnum1<0 9\else 1\fi\relax T\else F\fi}
\M{o=\unless\ifnum 5<\ifx\relax\relax 9\else 1\fi\relax T\else F\fi}
\M{p=\unless\ifnum 5<\ifx\relax\undefined 9\else 1\fi\relax T\else F\fi}
% two levels deep, both prefixed by \unless
\M{q=\unless\ifnum 2<\unless\ifodd 3 9\else 1\fi\relax T\else F\fi}
\bye
